\section{Reverse a Stack using Recursion}
\label{sec:rec-reverse-stack}

\problemstatement{Given a stack, reverse its order using \textbf{only}
the stack's own \textit{push} and \textit{pop} operations and
\textbf{recursion} -- no second stack, array, or queue is allowed as
auxiliary storage.}

\begin{patternbox}
\emph{``Solve using only recursion, no auxiliary data structure''} is a
direct instruction to use the \textbf{call stack itself} as the only
permitted extra storage: every value temporarily removed from the stack
must be held as a local variable on some recursive call's stack frame,
not in any explicit container. This requires a second, helper recursive
function -- inserting a value back at the \emph{bottom} of a stack is
itself a small recursive sub-problem, since a stack has no native
``insert at the bottom'' operation.
\end{patternbox}

\subsection*{Understanding the problem carefully}

The overall idea: pop the top element off, recursively reverse
\emph{everything else} below it, and then place the popped element back
-- but at the \emph{bottom} of the now-reversed remainder, since
reversal means the original top must end up at the very bottom of the
final stack. Placing a value at the bottom of a stack, using only
\textit{push}/\textit{pop}, is itself solved recursively: pop everything
off one at a time, push the target value once the stack is empty, then
push everything else back on top, in the reverse of the order it was
removed (which recursion handles automatically via the call stack's own
LIFO unwinding).

\begin{ideabox}[Two Cooperating Recursive Functions]
$\textsc{InsertAtBottom}(\textit{stack}, x)$: if $\textit{stack}$ is
empty, push $x$ and return (base case). Otherwise, pop the top into a
local variable $\textit{temp}$, recursively call
$\textsc{InsertAtBottom}(\textit{stack}, x)$ (placing $x$ at the bottom
of what remains), then push $\textit{temp}$ back on top.
$\textsc{ReverseStack}(\textit{stack})$: if $\textit{stack}$ is empty (or
has one element), return (base case). Otherwise, pop the top into
$\textit{temp}$, recursively call
$\textsc{ReverseStack}(\textit{stack})$ (reversing everything below
$\textit{temp}$), then call
$\textsc{InsertAtBottom}(\textit{stack}, \textit{temp})$.
\end{ideabox}

\subsection*{Why local variables on the call stack correctly substitute for an explicit auxiliary structure}

Each recursive call to $\textsc{ReverseStack}$ holds exactly one
popped element in its own local variable $\textit{temp}$, alive for the
duration of that call (including everything that happens inside its
recursive sub-call). Because the call stack itself maintains these
local variables in a strict LIFO order matching the order elements were
popped, and because each element is pushed back (via
$\textsc{InsertAtBottom}$) exactly once as its corresponding call
returns, no information is lost and no second container is ever needed
-- the recursion's own bookkeeping does the job an explicit auxiliary
stack would otherwise have done.

\subsection*{Algorithm}

\begin{enumerate}
  \item $\textsc{InsertAtBottom}(\textit{stack}, x)$:
  \begin{enumerate}
    \item If $\textit{stack}$ is empty: push $x$; return.
    \item Pop top into $\textit{temp}$.
    \item $\textsc{InsertAtBottom}(\textit{stack}, x)$.
    \item Push $\textit{temp}$.
  \end{enumerate}
  \item $\textsc{ReverseStack}(\textit{stack})$:
  \begin{enumerate}
    \item If $\textit{stack}$ is empty: return.
    \item Pop top into $\textit{temp}$.
    \item $\textsc{ReverseStack}(\textit{stack})$.
    \item $\textsc{InsertAtBottom}(\textit{stack}, \textit{temp})$.
  \end{enumerate}
\end{enumerate}

\subsection*{Dry run}

\dryrun{Reverse the stack $[1,2,3]$ (top to bottom: $3,2,1$).}

\begin{itemize}
  \item $\textsc{ReverseStack}$ call 1: pop $3$. Recurse on $[1,2]$.
  \item Call 2: pop $2$. Recurse on $[1]$.
  \item Call 3: pop $1$. Recurse on $[]$ (empty, base case, returns
        immediately).
  \item Back in call 3: $\textsc{InsertAtBottom}([], 1)$: stack is
        empty, push $1$. Stack $=[1]$.
  \item Back in call 2: $\textsc{InsertAtBottom}([1], 2)$: pop $1$ into
        \textit{temp}; recurse $\textsc{InsertAtBottom}([], 2)$: empty,
        push $2$, stack$=[2]$; back in outer call, push \textit{temp}
        ($1$) back: stack$=[2,1]$ (top to bottom: $1,2$).
  \item Back in call 1: $\textsc{InsertAtBottom}([2,1], 3)$: pop $1$
        into \textit{temp}; recurse on $[2]$: pop $2$ into another
        \textit{temp}; recurse on $[]$: push $3$, stack$=[3]$; push
        back the inner \textit{temp} ($2$): stack$=[3,2]$; push back
        the outer \textit{temp} ($1$): stack$=[3,2,1]$.
\end{itemize}

Final stack (top to bottom): $1,2,3$ -- the reverse of the original
$3,2,1$.

\subsection*{C++ implementation}

\begin{lstlisting}
#include <bits/stdc++.h>
using namespace std;

void insertAtBottom(stack<int>& st, int x) {
    if (st.empty()) {
        st.push(x);
        return;
    }
    int temp = st.top();
    st.pop();
    insertAtBottom(st, x);
    st.push(temp);
}

void reverseStack(stack<int>& st) {
    if (st.empty()) return;

    int temp = st.top();
    st.pop();
    reverseStack(st);
    insertAtBottom(st, temp);
}

int main() {
    stack<int> st;
    st.push(1); st.push(2); st.push(3);

    reverseStack(st);

    while (!st.empty()) {
        cout << st.top() << " ";
        st.pop();
    }
    cout << endl;
    return 0;
}
\end{lstlisting}

\begin{complexitybox}
\textbf{Time Complexity.} $\textsc{InsertAtBottom}$ costs $O(k)$ for a
stack of size $k$ (it must pop and re-push every element to reach the
bottom). $\textsc{ReverseStack}$ calls it once for each of the $n$
elements, with the stack shrinking by one each time, giving
\[
  O(1+2+\cdots+n) = O(n^2).
\]
\textbf{Space Complexity.} $O(n)$ for the recursion stack at the deepest
point (both functions' call depths are bounded by $n$).
\end{complexitybox}

\begin{takeawaybox}
``Using only recursion, no auxiliary data structure'' is a direct
instruction to use \emph{local variables on the call stack} as your only
permitted temporary storage. Many such ``stack-only'' problems
(reversing, sorting, this section and the next) need exactly this kind of
helper function that performs an operation a real stack does not natively
support (here, inserting at the bottom), built entirely out of recursive
pop/push/recurse/push-back steps.
\end{takeawaybox}
